\section{Related Work}
\label{sec:related_work}
% \textbf{Question answering over knowledge graphs.}
% simple QA is solved
The most commonly used KGQA benchmark is the SimpleQuestions \cite{DBLP:journals/corr/BordesUCW15} dataset, which contains questions that require identifying a single triple to retrieve the correct answers.
Recent results~\cite{DBLP:conf/emnlp/PetrochukZ18} show that most of these simple questions can be solved using a standard neural network architectures.
This architecture consists of two components: (1) a conditional random fields (CRF) tagger with GloVe word embeddings for subject recognition given the text of the question, and (2) a bidirectional LSTM with FastText word embeddings for relation classification given the text of the question and the subject from the previous component.
Approaches to simple KGQA cannot easily be adapted to solving complex questions, since they rely heavily on the assumption that each question refers to only one entity and one relation in the KG, which is no longer the case for complex questions.
Moreover, complex KGQA also requires matching more complex graph patterns beyond a single triple.

Since developing KGQA systems requires solving several tasks, namely entity, relation and class linking, and afterwards query building, they are often implemented as independent components and arranged into a single pipeline~\cite{DBLP:conf/semweb/DubeyBCL18}.
Frameworks such as QALL-ME~\cite{DBLP:journals/ws/FerrandezSKDFNITONMG11}, OKBQA~\cite{DBLP:conf/semweb/KimUNFHKCKUKC17} and Frankenstein~\cite{DBLP:conf/esws/SinghBRS18}, allow one to share and reuse those components as a collaborative effort.
For example, Frankenstein includes 29 components that can be combined and interchanged~\cite{DBLP:conf/www/SinghRBSLUVKP0V18}.
However, the distribution of the number of components designed for each task is very unbalanced.
Most of the components in Frankenstein support entity and relation linking, 18 and 5 components respectively, while only two components perform query building~\cite{DBLP:journals/corr/abs-1809-10044}.

There is a lack of diversity in approaches that are being considered for retrieving answers from a KG.
OKBQA and Frankenstein both propose to translate natural language questions to SPARQL queries and then use existing query processing mechanism to retrieve answers.\footnote{\url{http://doc.okbqa.org/query-generation-module/v1/}}
We show that using matrix algebra approaches is more efficient in case of natural language processing than traditional SPARQL-based approaches since they are optimized for parallel computation, thereby allowing us to explore multiple alternative question interpretations at the same time~\cite{DBLP:conf/hpec/KepnerABBFGHKLM16,jamour2018demonstration}.

Query building approaches involve query generation and ranking steps~\cite{maheshwari2018learning,DBLP:conf/esws/ZafarNL18}.
These approaches essentially consider KGQA as a subgraph matching task~\cite{DBLP:conf/coling/BaoDYZZ16,maheshwari2018learning,DBLP:conf/coling/SorokinG18}, which is an NP-hard problem (by reduction to the subgraph isomorphism problem)~\cite{DBLP:conf/sigmod/ZouHWYHZ14}.
In practice, \citet{DBLP:journals/corr/abs-1809-10044} report that the question building components of Frankenstein fail to process 46\% questions from a subset of LC-QuAD due to the large number of triple patterns.
The reason is that most approaches to query generation are template-based~\cite{DBLP:journals/corr/abs-1803-00832} and complex questions require a large number of candidate templates~\cite{DBLP:journals/corr/abs-1809-10044}.
For example, WDAqua~\cite{DBLP:journals/corr/abs-1803-00832} generates 395  SPARQL queries as possible interpretations for the question ``Give me philosophers born in Saint Etienne.''

In summary, we identify the query building component as the main bottleneck for the development of KGQA systems and propose \OurApproach{} as an alternative to the query building approach.
Also, the pipeline paradigm is inefficient since it requires KG access first for disambiguation and then again for query building.
\OurApproach{} accesses the KG only to aggregate the confidence scores via graph traversal after question parsing and shallow linking that matches an input question to labels of nodes and edges in the KG.

The work most similar to ours is the spreading activation model of Treo~\cite{DBLP:journals/dke/FreitasOOSC13}, which is also a no-SPARQL approach based on graph traversal that propagates relatedness scores for ranking nodes with a cut-off threshold.
Treo relies on POS tags, the Stanford dependency parser, Wikipedia links and TF/IDF vectors for computing semantic relatedness scores between a question and terms in the KG.
Despite good performance on the QALD 2011 dataset, the main limitation of Treo is an average query execution time of 203s~\cite{DBLP:journals/dke/FreitasOOSC13}.
In this paper we show how to scale this kind of approach to large KGs and complement it with machine learning approaches for question parsing and word embeddings for semantic expansion.

Our approach overcomes the limitations of the previously proposed graph-based approach in terms of efficiency and scalability, which we demonstrate on a compelling benchmark.
We evaluate \OurApproach{} on LC-QuAD~\cite{DBLP:conf/semweb/TrivediMDL17}, which is the largest dataset used for benchmarking complex KGQA.
WDAqua is our baseline approach, which is the state-of-the-art in KGQA as the winner of the most recent Scalable Question Answering Challenge (SQA2018)~\cite{DBLP:conf/esws/NapolitanoUN18}.
Our evaluation results demonstrate improvements in precision and recall, while reducing average execution time over the SPARQL-based WDAqua, which is also orders of magnitude faster than results reported for the previous graph-based approach Treo.

There is other work on KGQA that uses embedding queries into a vector space~\cite{DBLP:conf/nips/HamiltonBZJL18,DBLP:conf/semweb/WangWLCZQ18}.
The benefit of our graph-based approach is in preserving the original structure of the KG that can be used for executing precise formal queries and answering ambiguous natural language questions.
The graph structure also makes the results traceable and, therefore, interpretable in terms of relevant paths and subgraphs in comparison with vector space operations.

\OurApproach{} uses message passing, a family of approaches that were initially developed in the context of probabilistic graphical models~\cite{pearl1988probabilistic,koller2009probabilistic}.
Graph neural networks trained to learn patterns of message passing have recently shown to be effective on a variety of tasks~\cite{DBLP:conf/icml/GilmerSRVD17,DBLP:journals/corr/abs-1806-01261}, including KG completion~\cite{DBLP:conf/esws/SchlichtkrullKB18}.
We show that our unsupervised message passing approach performs well on complex question answering and helps to overcome sampling biases in the training data, which supervised approaches are prone to.
